\sectionkicker{Source-backed technical notes}
\subsection*{年次総括――2022年に変わった設計空間、変わらなかった制約: Technical Notes}
本欄は記事本文で圧縮した一次資料上の情報を、比較・再検証しやすい形で整理したものである。時系列、確認済みの主張、留意点、一次資料URLを掲載する。
\smallskip
\noindent{\small\textit{Technical Notesでは、一次資料から確認した公開・提供・研究上の事実を記載する。数値・能力評価や提供元・著者による比較表現は、独立再現済みの結果ではなく帰属claimとして扱う。}}\par
\medskip
\subsection*{Theme at a glance}
\addcontentsline{toc}{subsection}{Theme at a glance}
\begin{center}
\footnotesize
\begin{tabularx}{\linewidth}{@{}p{0.20\linewidth}p{0.10\linewidth}p{0.14\linewidth}X@{}}
\toprule
資料 & 位置づけ & 種別 & 時系列 \\
\midrule
InstructGPT & 補足資料 & モデル & 2022-01-27, 2022-03-04 \\
Chinchilla & 補足資料 & モデル & 2022-03-29, 2022-04-12 \\
DALL-E 2 & 補足資料 & モデル & 2022-04-13, 2022-07-20, 2022-09-28, 2022-11-03 \\
Stable Diffusion & 補足資料 & オープンウェイト & 2022-08-10, 2022-08-22 \\
ReAct & 補足資料 & 論文 & 2022-10-06 \\
ChatGPT research preview & 補足資料 & 製品 & 2022-11-30 \\
HELM & 補足資料 & 評価ベンチマーク & 2022-11-16 \\
\bottomrule
\end{tabularx}
\end{center}
\normalsize

\begin{technicalnote}{InstructGPT}{補足資料}
\begingroup
% reader-facing Technical Notes break policy
\widowpenalty=10000
\clubpenalty=10000
\displaywidowpenalty=10000
\begin{tabularx}{\linewidth}{@{}>{\bfseries}p{0.22\linewidth}X@{}}
組織 & OpenAI \\
種別 & モデル \\
時系列 & 2022-01-27, 2022-03-04 \\
\end{tabularx}
\smallskip
{\bfseries 一次資料から整理したtechnical points}
\begin{itemize}[leftmargin=1.5em,itemsep=0.35em]
\item \textbf{一次情報で確認できる事実}: OpenAIの選定済み一次資料では、labelerによる望ましい応答のdemonstrationでGPT-3をsupervised fine-tuningした後、model出力のranking dataをhuman feedbackとして用い、reinforcement learningでinstruction-following modelを調整する手順が記述されている。
\end{itemize}
\begin{minipage}{\linewidth}
% reader-facing Technical Notes generic boundary/source tail group
\begin{samepage}
{\bfseries 一次資料}
\begingroup\sloppy
\Urlmuskip=0mu plus 2mu\relax
\begin{itemize}[leftmargin=1.5em,itemsep=0.25em]
\item {\scriptsize\url{http://arxiv.org/abs/2203.02155v1}}
\item {\scriptsize\url{https://openai.com/index/instruction-following}}
\end{itemize}
\endgroup
\end{samepage}
\end{minipage}
\endgroup
\end{technicalnote}
\medskip

\begin{technicalnote}{Chinchilla}{補足資料}
\begingroup
% reader-facing Technical Notes break policy
\widowpenalty=10000
\clubpenalty=10000
\displaywidowpenalty=10000
\begin{tabularx}{\linewidth}{@{}>{\bfseries}p{0.22\linewidth}X@{}}
組織 & Google DeepMind \\
種別 & モデル \\
時系列 & 2022-03-29, 2022-04-12 \\
\end{tabularx}
\smallskip
{\bfseries 一次資料から整理したtechnical points}
\begin{itemize}[leftmargin=1.5em,itemsep=0.35em]
\item \textbf{一次情報で確認できる事実}: Google DeepMindの選定済み一次資料では、所与のtraining compute budgetの下でTransformer language modelのmodel sizeとtraining token数の配分を調べ、compute-optimal trainingでは両者を合わせて増やす必要があるというscaling relationを実験的に検討した。
\end{itemize}
\begin{minipage}{\linewidth}
% reader-facing Technical Notes generic boundary/source tail group
\begin{samepage}
{\bfseries 一次資料}
\begingroup\sloppy
\Urlmuskip=0mu plus 2mu\relax
\begin{itemize}[leftmargin=1.5em,itemsep=0.25em]
\item {\scriptsize\url{http://arxiv.org/abs/2203.15556v1}}
\item {\scriptsize\url{https://deepmind.google/blog/an-empirical-analysis-of-compute-optimal-large-language-model-training/}}
\end{itemize}
\endgroup
\end{samepage}
\end{minipage}
\endgroup
\end{technicalnote}
\medskip

\begin{technicalnote}{DALL-E 2}{補足資料}
\begingroup
% reader-facing Technical Notes break policy
\widowpenalty=10000
\clubpenalty=10000
\displaywidowpenalty=10000
\begin{tabularx}{\linewidth}{@{}>{\bfseries}p{0.22\linewidth}X@{}}
組織 & OpenAI \\
種別 & モデル \\
時系列 & 2022-04-13, 2022-07-20, 2022-09-28, 2022-11-03 \\
\end{tabularx}
\smallskip
{\bfseries 一次資料から整理したtechnical points}
\begin{itemize}[leftmargin=1.5em,itemsep=0.35em]
\item \textbf{一次情報で確認できる事実}: OpenAIの選定済み一次資料では、text captionからCLIP image embeddingを生成するpriorと、そのembeddingを条件に画像を生成するdecoderからなる二段構成を提示し、decoderにはdiffusion modelを用いた。2022年中のbeta、waitlist撤廃、API public betaは、この研究modelの提供範囲が段階的に拡張した別々のlifecycle eventとして記録する。
\end{itemize}
\begin{minipage}{\linewidth}
% reader-facing Technical Notes generic boundary/source tail group
\begin{samepage}
{\bfseries 一次資料}
\begingroup\sloppy
\Urlmuskip=0mu plus 2mu\relax
\begin{itemize}[leftmargin=1.5em,itemsep=0.25em]
\item {\scriptsize\url{http://arxiv.org/abs/2204.06125v1}}
\item {\scriptsize\url{https://openai.com/index/dall-e-api-now-available-in-public-beta}}
\item {\scriptsize\url{https://openai.com/index/dall-e-now-available-in-beta}}
\item {\scriptsize\url{https://openai.com/index/dall-e-now-available-without-waitlist}}
\item {\scriptsize\url{https://openai.com/index/hierarchical-text-conditional-image-generation-with-clip-latents}}
\end{itemize}
\endgroup
\end{samepage}
\end{minipage}
\endgroup
\end{technicalnote}
\medskip

\begin{technicalnote}{Stable Diffusion}{補足資料}
\begingroup
% reader-facing Technical Notes break policy
\widowpenalty=10000
\clubpenalty=10000
\displaywidowpenalty=10000
\begin{tabularx}{\linewidth}{@{}>{\bfseries}p{0.22\linewidth}X@{}}
組織 & Stability AI / collaborators \\
種別 & オープンウェイト \\
時系列 & 2022-08-10, 2022-08-22 \\
\end{tabularx}
\smallskip
{\bfseries 一次資料から整理したtechnical points}
\begin{itemize}[leftmargin=1.5em,itemsep=0.35em]
\item \textbf{一次情報で確認できる事実}: Stability AI / collaboratorsの選定済み一次資料では、Stability AIは2022年8月10日にStable Diffusionをまずresearcher向けにreleaseし、model weightsはaccess取得後にHugging Faceで提供、codeとmodel cardも公開する段階的releaseとして位置付けた。 8月22日のpublic releaseではweights・model card・codeを一般公開し、Creative ML OpenRAIL-M licenseの下でcommercial/non-commercial利用を認める公開境界へ移行した。
\end{itemize}
\begin{minipage}{\linewidth}
% reader-facing Technical Notes generic boundary/source tail group
\begin{samepage}
{\bfseries 一次資料}
\begingroup\sloppy
\Urlmuskip=0mu plus 2mu\relax
\begin{itemize}[leftmargin=1.5em,itemsep=0.25em]
\item {\scriptsize\url{https://stability.ai/news-updates/stable-diffusion-announcement}}
\item {\scriptsize\url{https://stability.ai/news-updates/stable-diffusion-public-release}}
\end{itemize}
\endgroup
\end{samepage}
\end{minipage}
\endgroup
\end{technicalnote}
\medskip

\begin{technicalnote}{ReAct}{補足資料}
\begingroup
% reader-facing Technical Notes break policy
\widowpenalty=10000
\clubpenalty=10000
\displaywidowpenalty=10000
\begin{tabularx}{\linewidth}{@{}>{\bfseries}p{0.22\linewidth}X@{}}
組織 & authors \\
種別 & 論文 \\
時系列 & 2022-10-06 \\
\end{tabularx}
\smallskip
{\bfseries 一次資料から整理したtechnical points}
\begin{itemize}[leftmargin=1.5em,itemsep=0.35em]
\item \textbf{一次情報で確認できる事実}: authorsの選定済み一次資料では、ReActはlanguage modelにreasoning traceとtask-specific actionをinterleaveして生成させ、推論側でaction planを追跡・更新しつつ、action側で外部knowledge baseやenvironmentから追加情報を取得する構成を提案した。 論文ではHotpotQA・FEVERでWikipedia APIとのinteractionを、ALFWorld・WebShopでinteractive decision makingを扱い、reasoningとactionを別々ではなく一つのtask-solving trajectoryとして接続した。
\end{itemize}
\begin{minipage}{\linewidth}
% reader-facing Technical Notes generic boundary/source tail group
\begin{samepage}
{\bfseries 一次資料}
\begingroup\sloppy
\Urlmuskip=0mu plus 2mu\relax
\begin{itemize}[leftmargin=1.5em,itemsep=0.25em]
\item {\scriptsize\url{http://arxiv.org/abs/2210.03629v3}}
\end{itemize}
\endgroup
\end{samepage}
\end{minipage}
\endgroup
\end{technicalnote}
\medskip

\begin{technicalnote}{ChatGPT research preview}{補足資料}
\begingroup
% reader-facing Technical Notes break policy
\widowpenalty=10000
\clubpenalty=10000
\displaywidowpenalty=10000
\begin{tabularx}{\linewidth}{@{}>{\bfseries}p{0.22\linewidth}X@{}}
組織 & OpenAI \\
種別 & 製品 \\
時系列 & 2022-11-30 \\
\end{tabularx}
\smallskip
{\bfseries 一次資料から整理したtechnical points}
\begin{itemize}[leftmargin=1.5em,itemsep=0.35em]
\item \textbf{一次情報で確認できる事実}: OpenAIの選定済み一次資料では、OpenAIは2022年11月30日にChatGPTをconversational interactionを前面に出したresearch previewとして公開し、follow-up questionへの応答、誤りの認識、誤った前提への異議、不適切なrequestの拒否をdialogue format上の挙動として説明した。 この公開は後年の一般的なChatGPT製品状態を遡及させるものではなく、2022年末時点の対話UIとresearch-preview availabilityの境界として扱う。
\end{itemize}
\begin{minipage}{\linewidth}
% reader-facing Technical Notes generic boundary/source tail group
\begin{samepage}
{\bfseries 一次資料}
\begingroup\sloppy
\Urlmuskip=0mu plus 2mu\relax
\begin{itemize}[leftmargin=1.5em,itemsep=0.25em]
\item {\scriptsize\url{https://openai.com/index/chatgpt}}
\end{itemize}
\endgroup
\end{samepage}
\end{minipage}
\endgroup
\end{technicalnote}
\medskip

\begin{technicalnote}{HELM}{補足資料}
\begingroup
% reader-facing Technical Notes break policy
\widowpenalty=10000
\clubpenalty=10000
\displaywidowpenalty=10000
\begin{tabularx}{\linewidth}{@{}>{\bfseries}p{0.22\linewidth}X@{}}
組織 & Stanford CRFM / authors \\
種別 & 評価ベンチマーク \\
時系列 & 2022-11-16 \\
\end{tabularx}
\smallskip
{\bfseries 一次資料から整理したtechnical points}
\begin{itemize}[leftmargin=1.5em,itemsep=0.35em]
\item \textbf{一次情報で確認できる事実}: Stanford CRFM / authorsの選定済み一次資料では、HELMはlanguage modelを少数のaccuracy指標だけで比較せず、複数scenarioと複数metricを組み合わせてholisticに評価するframeworkとして提案された。 評価面にはaccuracyだけでなくcalibration・robustness・fairness・bias・toxicity・efficiencyなどを含め、model comparisonのcoverageそのものをoperating constraintとして扱う。
\end{itemize}
\begin{minipage}{\linewidth}
% reader-facing Technical Notes generic boundary/source tail group
\begin{samepage}
{\bfseries 一次資料}
\begingroup\sloppy
\Urlmuskip=0mu plus 2mu\relax
\begin{itemize}[leftmargin=1.5em,itemsep=0.25em]
\item {\scriptsize\url{http://arxiv.org/abs/2211.09110v2}}
\end{itemize}
\endgroup
\end{samepage}
\end{minipage}
\endgroup
\end{technicalnote}
\medskip
